\subsection{Integration of functions with values in a Banach space}

THEOREM 1 (Lebesgue--Fubini). --- Let f be a function defined on T × T', with values in a Banach space F or in \(\overline{R}\) ; let N be the set of \(t \in T\) such that the function \(t' \mapsto \mathbf{f}(t, t')\) is not \(\mu'\) -integrable.

\begin{enumerate}
\def\labelenumi{\alph{enumi})}
\tightlist
\item
  Suppose that \(\mathbf{f}\) is \(\nu\) -integrable; then \(\mathrm{N}\) is \(\mu\) -negligible, the function \(t \mapsto \int \mathbf{f}(t, t') d\mu'(t')\) (defined for \(t \notin \mathrm{N}\) ) is \(\mu\) -integrable, and
\end{enumerate}

\[
\iint \mathbf {f} (t, t ^ {\prime}) d \mu (t) d \mu^ {\prime} (t ^ {\prime}) = \int d \mu (t) \int \mathbf {f} (t, t ^ {\prime}) d \mu^ {\prime} (t ^ {\prime}).\tag{13}
\]

\begin{enumerate}
\def\labelenumi{\alph{enumi})}
\setcounter{enumi}{1}
\tightlist
\item
  Suppose that \(\mathbf{f}\) is essentially \(\nu\) -integrable, and that the measure \(\mu'\) is moderated; then \(\mathbf{N}\) is locally \(\mu\) -negligible, the function \(t \mapsto \int \mathbf{f}(t, t') d\mu'(t')\) (defined for \(t \notin \mathbf{N}\) ) is essentially \(\mu\) -integrable, and (13) holds.
\end{enumerate}

The assertion a) follows at once from Th. 1 of §3, No.~3. To establish b), let us denote by g a \(\nu\) -integrable function equal to f locally almost everywhere, and by H the set of \((t, t')\) such that \(\mathbf{f}(t, t') \neq \mathbf{g}(t, t')\) . By Cor. 1 of Prop. 7, the section \(\mathrm{H}(t)\) is \(\mu'\) -negligible, except for \(t \in T\) forming a locally \(\mu\) -negligible set. The result pertaining to f may then be deduced from the statement a) applied to g.

Scholium. --- Let f be a function defined on \(T \times T'\) , with values in \(\overline{R}\) or in a Banach space, that is \(\nu\) -measurable and \(\nu\) -moderated. For the three integrals

\[
\iint \mathbf {f} (t, t ^ {\prime}) d \mu (t) d \mu^ {\prime} (t ^ {\prime}), \int d \mu (t) \int \mathbf {f} (t, t ^ {\prime}) d \mu^ {\prime} (t ^ {\prime}), \int d \mu^ {\prime} (t ^ {\prime}) \int \mathbf {f} (t, t ^ {\prime}) d \mu (t)
\]

to exist and be equal, it is necessary and sufficient that one of the two numbers

\[
\int^ {*} d \mu (t) \int^ {*} | {\bf f} (t, t ^ {\prime}) | d \mu^ {\prime} (t ^ {\prime}), \quad \int^ {*} d \mu^ {\prime} (t ^ {\prime}) \int^ {*} | {\bf f} (t, t ^ {\prime}) | d \mu (t)
\]

be finite.

This is an immediate consequence of Th. 1, Prop. 7 and the integrability criterion (Ch. IV, §5, No.~6, Th. 5).

Remarks. --- 1) When the measure \(\mu'\) is not moderated, it can happen that \(\mathbf{f}\) is essentially \(\nu\) -integrable and the function \(t' \mapsto \mathbf{f}(t, t')\) is not essentially \(\mu'\) -integrable for any value of \(t \in \mathrm{T}\) (§3, Exer. 4).

\begin{enumerate}
\def\labelenumi{\arabic{enumi})}
\setcounter{enumi}{1}
\tightlist
\item
  Let \(\mu\) and \(\mu'\) be two complex measures, and let \(\nu = \mu \otimes \mu'\) . If \(\mathbf{f}\) is \(\nu\) -integrable (in other words, \(|\nu|\) -integrable), then application of the theorem to the measures \(|\mu|\) and \(|\mu'|\) , whose product is \(|\nu|\) (Ch. III, §4, No.~2, Prop. 3), implies that \(t' \mapsto \mathbf{f}(t, t')\) is \(\mu'\) -integrable for \(\mu\) -almost every \(t\) . From this one deduces, on decomposing the measures \(\mu\) and \(\mu'\) as a linear combination of positive measures, that the statement of a) extends to complex measures. One can argue similarly for b).
\end{enumerate}

PROPOSITION 9. --- Let F, F' and G be three Banach spaces, and let \((\mathbf{x},\mathbf{y})\mapsto[\mathbf{x}\cdot\mathbf{y}]\) be a continuous bilinear mapping of \(F\times F'\) into G. Let f (resp. \(f'\) ) be a function defined on T (resp. \(T'\) ) with values in F (resp. \(F'\) ) and essentially integrable for \(\mu\) (resp. \(\mu'\) ). Let g be the function \((t,t')\mapsto[\mathbf{f}(t)\cdot\mathbf{f}'(t')]\) ; then g is essentially integrable for \(\mu\otimes\mu'\) , and

\[
\iint \left[ \mathbf {f} (t) \cdot \mathbf {f} ^ {\prime} \left(t ^ {\prime}\right) \right] d \mu (t) d \mu^ {\prime} \left(t ^ {\prime}\right) = \left[ \left(\int \mathbf {f} d \mu (t)\right) \cdot \left(\int \mathbf {f} ^ {\prime} \left(t ^ {\prime}\right) d \mu^ {\prime} \left(t ^ {\prime}\right)\right) \right].\tag{14}
\]

If, moreover, \(\mathbf{f}\) and \(\mathbf{f}'\) are integrable, then \(\mathbf{g}\) is integrable.

The function \((t, t') \mapsto [\mathbf{f}(t) \cdot \mathbf{f}'(t')]\) is \((\mu \otimes \mu')\) -measurable by Cor. 1 of Prop. 3. On the other hand, if \(b\) denotes the norm of the bilinear mapping \((\mathbf{x}, \mathbf{y}) \mapsto [\mathbf{x} \cdot \mathbf{y}]\) , then

\[
\begin{array}{r l} \iint^ {\bullet} | [ \mathbf {f} (t) \cdot \mathbf {f} ^ {\prime} (t ^ {\prime}) ] | d \mu (t) d \mu^ {\prime} (t ^ {\prime}) & \leqslant b \iint^ {\bullet} | \mathbf {f} (t) | \cdot | \mathbf {f} ^ {\prime} (t ^ {\prime}) | d \mu (t) d \mu^ {\prime} (t ^ {\prime}) \\ & = b \left(\int^ {\bullet} | \mathbf {f} (t) | d \mu (t)\right) \left(\int^ {\bullet} | \mathbf {f} ^ {\prime} (t ^ {\prime}) | d \mu^ {\prime} (t ^ {\prime})\right) \end{array}
\]

by virtue of Prop. 8. This shows that \([\mathbf{f}(t) \cdot \mathbf{f}'(t')]\) is essentially integrable for \(\mu \otimes \mu'\) (§1, No.~3, Prop. 9). Suppose that \(\mathbf{f}\) and \(\mathbf{f}'\) are integrable: then \(\mathbf{f}\) and \(\mathbf{f}'\) are moderated, and \(\mathbf{g}\) is moderated (Cor. 2 of Prop. 5) therefore integrable (§1, No.~3, Cor. of Prop. 9). In this case the formula (14) follows from the Lebesgue--Fubini theorem and the linearity of the integral (Ch. IV, §4, No.~2, Th. 1). To complete the treatment of the case that \(\mathbf{f}\) and \(\mathbf{f}'\) are essentially integrable, one then applies (14) to two integrable functions \(\mathbf{f}_1\) and \(\mathbf{f}_1'\) , equal locally almost everywhere to \(\mathbf{f}\) and \(\mathbf{f}'\) , on observing that \([\mathbf{f} \cdot \mathbf{f}'] = [\mathbf{f}_1 \cdot \mathbf{f}_1']\) locally almost everywhere in \(T \times T'\) (Prop. 4).

This result extends to the product of complex measures.

